\subsection{Live Test}
CryptoGuard was also run end-to-end on a 2-vCPU Linux cloud container with active response (process termination on, network isolation off). A baseline of 120 windows of 5\,s was recorded (mean CPU 0.3\%), the model was trained on it, and the harmless miner simulator was then started twice. In the \emph{noisy} run (all cores, \texttt{xmrig} in the command line, connection to a fake pool on port 3333), the first alert came 0.8\,s after the miner started and the process was terminated after 0.8\,s. In the \emph{stealthy} run (one core at 50\%, no signature, no pool connection), detection relied on behaviour alone: the first alert came after 10.8\,s and the process was terminated after 91.5\,s, once the drift confidence passed the kill threshold. No alerts were raised in the windows before either miner started. Decisions are taken at the end of each window, so these latencies depend on when in a window the miner starts. The live test shows that the pipeline works end to end on a real host; it is not a measurement of accuracy, because the idle test host is far simpler than a production server.

The first alert of the stealthy run, raised from behaviour alone, read: \emph{Looks like cryptojacking because the server is using 25 percentage points more CPU than its current traffic and process activity can explain (normally under 1); CPU usage is 26\%, while this server normally runs at about 0\% at 15:00; a single process is using 25\% of total CPU (normal top process: ~0\%).}
